\documentclass[10pt,a4paper]{amsart}
\usepackage[margin=19mm]{geometry}
\usepackage{amsmath,amssymb,mathtools}
\usepackage[scr=rsfs,cal=euler]{mathalfa}
\usepackage{xfrac}
\usepackage[unicode,hidelinks,bookmarks=false]{hyperref}
\newcommand{\ie}{i.e.\ }
\newcommand{\cf}{cf.\ }
\newcommand{\resp}{respectively\ }
\newcommand{\Subsection}{\subsection}
\providecommand{\nobreakdash}{\nobreak\hbox{-}}
\setlength{\emergencystretch}{2em}
\multlinegap0pt
\usepackage{mathtools}
\usepackage{amssymb}
\usepackage[normalem]{ulem}
\newtheorem{proposition}{Proposition}
\DeclareMathOperator{\im}{im}
\title{Computing Projective Implicit Representations from Poset Towers}
\author{Tamal K. Dey, Florian Russold}
\date{Source: arXiv:2505.08755v2 (CC BY 4.0)}
\begin{document}
\maketitle

\noindent\emph{Source and licence.} Computing Projective Implicit Representations from Poset Towers. Version arXiv:2505.08755v2 (CC BY 4.0), distributed by arXiv under the Creative Commons Attribution 4.0 International licence, http://creativecommons.org/licenses/by/4.0/, as recorded in the arXiv licence metadata for this exact version and shown on the abstract page https://arxiv.org/abs/2505.08755v2. Full licence notice: \texttt{licenses/arxiv-2505.08755-CC-BY-4.0.txt}.\par\smallskip

\noindent\textit{Editorial scope.} Counted unit: the article's proposition prop:graph\_equation (source0.tex lines 1539--1541). The article's complete original proof of it is delivered with it (source0.tex lines 1543--1593, one \texttt{proof} environment of 51 lines). No mathematics is written, repaired or re-derived: the statement and the proof are the article's own lines, in the article's own notation, and no retained line is substituted or re-flowed.  No further source-local context is needed: every cross-reference of every retained line resolves inside the two delivered spans.  Nothing was omitted from any retained span, so the package carries no omission record.  No retained line contains a citation, so the wrapper carries no bibliography and no bibliography entry.  No retained line contains a figure environment, an image-inclusion command or drawing code, so no image is delivered and the package declares no asset dependency.  Editorial formatting only, in addition: an AMS \texttt{amsart} preamble that restates the article's own declarations and macros used by the retained lines, namely the environments \texttt{proposition} on separate counters, together with the standard packages those lines need.  The retained environments print in this document as Proposition 1 in order of appearance, whereas the article prints the same items with the numbering of its own full document.  The complete native source of the article as distributed in the arXiv e-print for this exact version is bundled unchanged as \texttt{sources/178064/source0.tex}.
\begin{proposition} \label{prop:graph_equation}
Let $Ax=b$ be a system of linear equations such that $A$ is the boundary matrix of a graph $G$ (without self-loops) and $b\in \im(A)$. Let $F=\bigcup_{i=1}^n F^i$ be a spanning forest of $G$ with connected components $F^1,\ldots,F^n$. Moreover, let $A^i$ be the submatrix of $A$ corresponding to the edges and vertices in $F^i$, $x^i$ the subvector of variables $x$ corresponding to the edges in $F^i$, $x^r$ the subvector of $x$ corresponding to edges in $G\setminus F$, and $b^i$ the subvector of $b$ corresponding to the vertices in $F^i$. Then $b^i\in\im(A^i)$ and $A^ix^i=b^i$ and $x^r=0$ implies $Ax=b$.  
\end{proposition}

\begin{proof}
Let $G^1,\ldots,G^n$ be the connected components of the graph $G$. Let $\hat{A}^i$ be the submatrix of $A$ corresponding to the vertices and edges in $G^i$, $\hat{x}^i$ the subvector of variables $x$ corresponding to the edges in $G^i$, and $\hat{b}^i$ the subvector of $b$ corresponding to the vertices in $G^i$. After reordering the columns and rows of $A$, $x$ and $b$ we can write the system as
\begin{equation*}
Ax=
\begin{pmatrix}
\hat{A}^1 & \cdots & 0 \\
\vdots & \ddots & \vdots \\
0 & \cdots & \hat{A}^n
\end{pmatrix}
\begin{pmatrix}
\hat{x}^1 \\
\vdots \\
\hat{x}^n \\
\end{pmatrix}=
\begin{pmatrix}
\hat{A}^1\hat{x}^1 \\
\vdots \\
\hat{A}^n\hat{x}^n \\
\end{pmatrix}=
\begin{pmatrix}
\hat{b}^1 \\
\vdots \\
\hat{b}^n 
\end{pmatrix}=b ,
\end{equation*}
where the coefficient matrix is in block-diagonal form because the connected components $G^1,\ldots,G^n$ of $G$ partition the vertices and edges into a disjoint union. Therefore, $Ax=b$ has a solution if and only if $\hat{A}^i\hat{x}^i=\hat{b}^i$ has a solution for all $1\leq i\leq n$. In other words, $b\in\im(A)$ if and only if $\hat{b}^i\in\im(\hat{A}^i)$. 

Since $F$ is a spanning forest of $G$, each connected component $F^i$ of $F$ must be a spanning tree of some connected component of $G$. W.l.o.g.\ assume $F^i$ is a spanning tree of the component $G^i$. Hence, $A^i$ is the submatrix of $\hat{A}^i$ corresponding to the edges in $F^i$. If $C^i$ is the submatrix of $\hat{A}^i$ corresponding to the edges in $G^i\setminus F^i$, then we can write $\hat{A}^i=\begin{pmatrix}A^i & C^i \end{pmatrix}$. Every edge in $G^i\setminus F^i$ closes a cycle in $G^i$ with the unique path connecting its endpoints in the spanning tree $F^i$. Thus, every column of $C^i$ is a linear combination of columns in $A^i$, given by the columns corresponding to the edges on this path. This implies that $\im(A^i)=\im(\hat{A}^i)$ and, since $G^i$ and $F^i$ have the same vertices, $b^i=\hat{b}^i\in \im(\hat{A}^i)=\im(A^i)$.  

After reordering the columns and rows of $A$, $x$ and $b$ again, we can write the system as
\begin{equation*}
Ax=
\begin{pmatrix}
A^1 & \cdots & 0 & C^1 & \cdots & 0 \\
\vdots & \ddots & \vdots & \vdots & \ddots & \vdots \\
0 & \cdots & A^n & 0 & \cdots & C^n
\end{pmatrix}
\begin{pmatrix}
x^1 \\
\vdots \\
x^n \\
x^r
\end{pmatrix}=
\begin{pmatrix}
b^1 \\
\vdots \\
b^n 
\end{pmatrix}=b ,
\end{equation*}
where $C^1,\ldots, C^n$ are the submatrices and $x^r$ is the subvector corresponding to all edges in $G\setminus F$. Since the columns of $A^1,\ldots,A^n$ correspond to edges in the connected components of the spanning forest $F$, the block of $A$ containing them is block-diagonal. Therefore, $A^1x^1=b^1,\ldots, A^nx^n=b^n$ and $x^r=0$ implies $Ax=b$.  
\end{proof}

\end{document}
